\documentclass[a4paper,11pt]{jsarticle}
\usepackage[utf8]{inputenc}
\usepackage{amssymb} % チェックボックス用

\begin{document}

{\Large \textbf{生成AI 利用申告書}}

\vspace{1em}

\begin{itemize}
    \item \textbf{氏名}: 梅澤颯太
    \item \textbf{学籍番号}: 25G1021
    \item \textbf{プログラムファイル名}: HCK02\_01.ino
    \item \textbf{使用した AI ツール}: Gemini, Claude
\end{itemize}

\section*{1. 何に使ったか}
今回の課題において，以下の目的で生成AIを利用しました．
\begin{itemize}
    \item シリアルモニタに正しく波形が表示されていない原因の特定．
    \item エラー修正
\end{itemize}

\section*{2. AIへの指示（プロンプト）}
AIに対して，主に以下のような質問・指示を出しました．
\begin{itemize}
    \item 「モニタに正しく波形が表示されていないらしいんだけどその原因は？」
    \item 「ちなみに出力を Serial.print(",振幅:"); Serial.print(a); だけにすると波形が出ないんだけどなんで？」
\end{itemize}

\section*{3. AI の出力の使い方}
AI の出力をどのように使いましたか．
\begin{itemize}
    \item[$\square$] そのまま使った
    \item[$\boxtimes$] 一部修正して使った
    \item[$\square$] 参考にしただけ
\end{itemize}
\textbf{理由}:
いくつか修正案を実行し，試したが変化のないものがあったので削除した．

\section*{4. 正しいと確認した方法}
AI の出力が正しいと，どのように確認しましたか．
\begin{itemize}
    \item 実際にarduinoに書き込み，実行して波形を確認した．
    \item 詳しくコードについてAIに質問したり，自分で調べることで理解を深めた．
\end{itemize}

\section*{5. 問題や間違い}
AIの出力に問題や間違いはありましたか．
\begin{itemize}
    \item[$\boxtimes$] あった
    \item[$\square$] なかった
\end{itemize}
\textbf{確認した内容や気づいたこと}:
量子化を10bitから14bitに変更するよう指示されたが，p.read関数を使っていて0-255までのデータしか使えないので必要がなかった．

\section*{6. 自分で工夫した点}
最終的なプログラムの中で，自分で考えて工夫したところを書いてください．
\begin{itemize}
    \item 必要な出力だけに絞って不必要な出力は消した．
\end{itemize}

\section*{参考文献}


\end{document}